\clearpage
\typeout{CITH-SUPPLEMENT-START}
\label{page:supplement-start}
\part*{Supplementary Material}
\suppressfloats[t]
\setcounter{section}{0}
\renewcommand{\thesection}{S\arabic{section}}
\renewcommand{\theHsection}{supp.\arabic{section}}
\renewcommand{\theHfigure}{supp.\arabic{figure}}
\renewcommand{\theHtable}{supp.\arabic{table}}
\setcounter{figure}{0}\renewcommand{\thefigure}{S\arabic{figure}}
\setcounter{table}{0}\renewcommand{\thetable}{S\arabic{table}}
\section{Retained workload and budget illustrations}
\label{supp:workload}
The numerical examples below are retained from v7.3. Their figures and frozen
arrays are unchanged; only their placement and notation have been updated.
Four controlled examples use $b=2$, $d(t)=2+\lfloor t/10\rfloor$ on $0\leq t\leq40$, and local capacities $(32,256,24)$ for update work, stored-state memory, and heat rejection. All $u_k=1$. Processing and thermal efficiency either grow as $e^{0.12t}$, grow as $e^{0.01t}$, or grow as $e^{0.12t}$ up to a ceiling of four; memory efficiency remains one. A fourth example activates at most eight worlds at a time without efficiency improvement. Every example stores all $W_d$ states. These are normalized accounting constructions, not device specifications.

The matched-efficiency example remains feasible while the full hierarchy grows from seven to 127 worlds. Demand outpacing efficiency first violates a constraint at $t=20$, and the efficiency-ceiling example first violates one at $t=40$. Scheduling remains feasible on the tested interval, but only eight worlds are active at once. The maximum ratios are respectively $0.4961$, $3.5471$, $1.3229$, and $0.4961$. All changes in depth occur at known times, so interval maxima can be checked at the jumps rather than inferred from a sampled time grid alone. \Cref{fig:support} illustrates the distinction between technology-supported expansion of internal activity and merely postponing a resource limit. Extending the matched-efficiency path with the same fixed memory efficiency and capacity gives $W=255$ and a memory ratio $255/256=0.99609375$ at $t=50$, but $W=511$ and $511/256=1.99609375$ at $t=60$. The latter is its first memory violation. The example therefore supports the stated finite horizon, not indefinite local sufficiency.

\begin{figure}[tbp]
\centering\includegraphics[width=.95\textwidth]{fig1.pdf}
\caption{Finite-horizon local support constructions. The largest service-demand ratio must remain at or below one. Processing and heat are charged to active worlds; persistent memory is charged to all represented states, including inactive ones. Growth in depth is stepwise. No curve is extrapolated to unlimited time, and no physical implementation is validated.}
\label{fig:support}
\end{figure}

\subsection{A joint finite-interval witness}

The support and allocation examples need not be merely disconnected illustrations. Couple the matched-efficiency schedule in \cref{sec:support} to $B=10$, $C_0=0$, $a(t)=W_{d(t)}/\ln(11)$, $r_n=r_h=0.1$, and all three prices equal to one. Set $q=0.99$, $\epsilon_R=\epsilon_I=0.001$, with unit initial slopes for $E,R,I$. At $z=B$, useful computation is exactly $C_L=W_{d(t)}$, the workload charged in the support example. The three outward-to-local marginal-return ratios are
\begin{equation}
 11\ln(11)\left[S_j+\frac{r_j+S_j}{W_{d(t)}}\right],
 \qquad(S_n,S_h,S_m)=(0.012,0.002,0.001).
 \label{eq:joint-witness}
\end{equation}
They decrease with world count and are already below one at seven worlds. Thus this same finite family satisfies both the three-service feasibility test and strict all-inward optimality with positive residual external values. As the hierarchy grows from seven to 127 worlds, attained logarithmic value rises from $\ln8$ to $\ln128$. The budget is fixed at $B=10$ and productivity is deliberately tuned to $a(t)=W_{d(t)}/\ln 11$. This is a consistency witness, not evidence for the plausibility of that productivity path, an engineering implementation, or an infinite-horizon optimum. Persistent memory is still charged to every world; future resource effects are considered separately in \cref{sec:dynamic}.

\subsection{Budget growth and shared technological progress}

With perfect experiential substitution and no residual external value, $q=1$ and $\epsilon_R=\epsilon_I=0$. For \emph{any} strictly increasing $F$, maximizing value is equivalent to maximizing $C$. Let
\begin{equation}
 R_* = \max\{r_n/p_n,r_h/p_h\}.
 \label{eq:Rstar}
\end{equation}
When $R_*>0$, the solution allocates
\begin{equation}
 z^*=\min\{B,\max(a/R_*-1,0)\},\qquad
 B_{\rm crit}=a/R_*-1.
 \label{eq:budgetcrit}
\end{equation}
Any remaining budget goes to a best-return computational external channel. The all-inward condition is $a/(1+B)\geq R_*$. If $R_*=0$, all investment is local. This exposes a limitation of a fixed-budget productivity path: with fixed $a$ and $R_*>0$, a sufficiently large budget restores external investment even without authenticity, information, or resilience premiums. The civilization does not stop preferring more value when it can generate virtual worlds.

If local and remote computational coefficients all receive the same multiplicative improvement, the pure-computation condition is unchanged. Progress that disproportionately raises intensive return can support inward allocation; progress that improves everything equally need not. \Cref{fig:budget} shows this budget effect for $R_*=1$. Persistent inward allocation along a development path requires the relevant inequalities to hold along that path, not merely at one chosen value of $B$.

\begin{figure}[tbp]
\centering\includegraphics[width=.95\textwidth]{fig2.pdf}
\caption{Exact inward budget share with perfect experiential substitution, zero residual value, and $R_*=1$. The transition occurs at $B=a-1$. Greater local productivity moves the boundary, but budget growth can eventually favor additional external capacity. Any strictly increasing computational value function gives the same allocation in this pure-computation case.}
\label{fig:budget}
\end{figure}

\section{Paired sensitivity design and numerical specification}
\label{supp:paired}\label{sec:access}
\subsection{Paired sensitivity design}
To distinguish deployment costs from access delays, use the normalized
family
\begin{equation}
 p_j=p_{j0}\left(f+\frac{1}{s_t}+\frac{0.2}{s_c}\right),\qquad
 r_j=r_{j0}m_r\exp(-0.1/s_t),\qquad a=a_0m_a.
 \label{eq:speed-map}
\end{equation}
Here $f>0$ is the residual price factor; $s_t,s_c>0$ index effective
material and information access speeds relative to $c$.
At instantaneous access, $p_j=p_{j0}f$ and $r_j=r_{j0}m_r$.
The coefficients $0.2$ and $0.1$ define scenarios, not calibrated physical
laws. Speed sweeps also remove the stipulated attenuation of remote
returns, so they are \emph{not} pure price experiments. The fixed-access
sweeps in $f$ isolate the price effect.

Faster-than-light access is only a counterfactual limiting case of this
map, not a physical mechanism established here. Instantaneous arrival
still leaves the positive deployment price $p_{j0}f$. Exactly zero
external prices are excluded: they can destroy compactness of the
feasible set and, with positive external computational returns and
unbounded increasing value, need not admit a finite optimizer.

The reference design contains 1,024 scrambled Sobol points in 14
dimensions, paired across every scenario. The exact library, generator
call and frozen inputs are specified in \cref{app:numerical}. Marginals
are log-uniform over the stated normalized ranges. The concave residual
benefits are
\begin{equation}
 E(x)=A_E(1-\exp(-x/s_E)),\quad
 R(u)=A_R(1-\exp(-u/s_R)),\quad
 I(v)=A_I(1-\exp(-v/s_I)).
 \label{eq:benefit-family}
\end{equation}
The reference uses $(q,\epsilon_R,\epsilon_I,f)=(0.99,0.01,0.05,1)$
and $m_a=m_r=1$. \Cref{tab:scenarios} retains the full paired counts.

At instantaneous access, the reference has 219 inward corners; a
uniform $10\%$ price cut leaves 207. Only 364 reference inputs admit
a finite productivity threshold with the same residual slopes; the other
660 violate a residual-floor condition. Perfect experiential substitution
with zero residual weights yields 527 inward corners, while 497 still
prefer external computational capacity. Raising only local productivity
tenfold gives 314 inward corners; raising local and remote productivity
together gives 222. At $f=0.01$ and $f=10^{-4}$, the count is zero.
These are exact classifications of the archived inputs, not estimates of
civilizational prevalence.

Eight additional scrambles of 4,096 points each give instantaneous
inward fractions of \DesignLogMin--\DesignLogMax\% under the same
log-uniform marginals. Linear-uniform marginals over the same endpoints
give \DesignLinearMin--\DesignLinearMax\%. These 65,536 classifications
use the exact boundary. Their variation describes sensitivity to the
parameter-space measure, not observational uncertainty about the Galaxy.

\subsection{Frozen design and retained paired results}
\label{supp:legacy-numerical}

The canonical unit points and transformed parameters are the arrays
\texttt{sobol\_vectors} and \texttt{parameters} in
\texttt{data/allocation/design.npz}. The generating environment records
Python~3.13.5, NumPy~2.3.5 and SciPy~1.17.0. The exact call is
\begin{verbatim}
from scipy.stats import qmc
u = qmc.Sobol(d=14, scramble=True,
              seed=20260905).random_base2(m=10)
\end{verbatim}
The default \texttt{bits=None} resolves to 30 bits; optimization is
\texttt{None}. No points are skipped or thinned. For the ordered endpoint
vectors in \cref{tab:ranges}, the transformation is
\begin{verbatim}
lo = np.log(lower_endpoints)
hi = np.log(upper_endpoints)
parameters = np.exp(lo + u * (hi - lo))
\end{verbatim}
Every scenario uses the same transformed rows.

Re-execution in SciPy~1.17.0 reproduces the archived unit points bit for
bit. Floating-point exponentiation changes some transformed entries by
at most a relative $2.22\times10^{-16}$ in the present environment;
all classifications agree. Evaluation of the frozen arrays reproduces
219 instantaneous-reference inward corners and 527 under perfect
substitution with zero residual weights, as well as all other counts
in \cref{tab:scenarios}.

With \texttt{seed=20260905}, the archived unit points are reproduced
bit for bit in the historical checks under SciPy~1.17.0 and 1.18.1;
\texttt{rng=20260905} yields a different scramble, with 208
instantaneous and 530 perfect-substitution inward corners.
The frozen arrays, not the seed alone, define the reference experiment.
The frozen package and the accompanying integrity manifest record the generator and checksums. The present rerun uses SciPy~1.17.0; the 1.18.1 result is retained archival evidence.

The parameter order is $(B,C_0,a,r_n,r_h,A_E,A_R,A_I,s_E,s_R,s_I,p_{n0},p_{h0},p_{m0})$. All base marginal ranges are log-uniform in the reference experiment. The stored columns named e, r, and i are written $s_E,s_R,s_I$ here; their values and ordering are unchanged.

\begin{table}[htbp]
\centering
\caption{Normalized ranges in the 14-dimensional reference design.}
\label{tab:ranges}
\begin{tabular}{@{}lll@{}}
\toprule
Quantity & Range & Meaning\\
\midrule
$B$ & $1$--$100$ & common investment budget\\
$C_0$ & $0.1$--$10$ & existing useful computational service\\
$a$ & $1$--$1000$ & intensive computational coefficient\\
$r_n,r_h$ & $0.1$--$10$ & external computational returns\\
$A_E,A_R,A_I$ & $0.1$--$10$ & benefit amplitudes\\
$s_E,s_R$ & $1$--$30$ & experiential and resilience scales\\
$s_I$ & $1$--$100$ & information scale\\
$p_{n0}$ & $1$--$10$ & resident unit price\\
$p_{h0}$ & $0.1$--$3$ & industry-only unit price\\
$p_{m0}$ & $0.01$--$0.3$ & observation-only unit price\\
\bottomrule
\end{tabular}
\end{table}

The sublight scenario has $(s_t,s_c)=(0.01,1)$, the communication-only variant $(0.01,\infty)$, the light-speed case $(1,1)$, and combined faster-access cases $(10,10)$, $(1000,1000)$, and $(\infty,\infty)$. The no-substitution variant sets $q=0$. Large residual values set $\epsilon_R=\epsilon_I=1$, either retaining $q=0.99$ or additionally setting $q=0$. The two tenfold productivity variants set $(m_a,m_r)=(10,1)$ and $(10,10)$. The perfect-substitution variant sets $(q,\epsilon_R,\epsilon_I)=(1,0,0)$. Reduced-price variants set $f=0.01$ or $10^{-4}$ with instantaneous access; other parameters remain at reference settings. A material-speed variable is not allowed to make numerical prices exactly zero.

\begin{table}[htbp]
\centering\small
\caption{All-inward inputs out of 1,024, with a conditional price-robustness check. The last column applies a uniform 10\% price reduction to each listed scenario while retaining its other settings. It is computed from the exact boundary, not from a new behavioral assumption or an observational sample.}
\label{tab:scenarios}

% BEGIN expanded scenario_table.tex
\begin{tabularx}{\textwidth}{@{}Xrr@{}}
\toprule
Scenario & All inward & Still inward after 10\% price cut \\
\midrule
Sublight material travel & 929 & 923 \\
Fast communication only & 929 & 923 \\
Light-speed limit & 326 & 313 \\
Travel and communication at 10c & 234 & 219 \\
Travel and communication at 1000c & 219 & 207 \\
Instantaneous reference & 219 & 207 \\
No experiential substitution & 146 & 135 \\
Large residual values & 14 & 14 \\
No substitution + large residuals & 11 & 10 \\
Local productivity $\times10$ & 314 & 296 \\
Local and remote productivity $\times10$ & 222 & 211 \\
Perfect substitution; no residual values & 527 & 505 \\
Outward price floor $\times0.01$ & 0 & 0 \\
Outward price floor $\times0.0001$ & 0 & 0 \\
\bottomrule
\end{tabularx}

% END expanded scenario_table.tex

\end{table}

Corner classification uses the exact inequalities, not optimizer-selected shares. The stored optimal shares are checked separately using the objective and concavity certificates described in the archived Supplementary Material and audit outputs. Small platform-dependent changes in transformed parameters or nearly tied optimized shares do not change any archived corner classification in the checks reported here.

For design-measure sensitivity, Sobol seeds 20260911 through 20260918 each generate $2^{12}$ points. Both log-uniform and linear-uniform transformations use the same bounds in \cref{tab:ranges}. Repeated scrambles assess numerical integration stability under a selected measure. They do not supply a prior justified by extraterrestrial observations.


\section{Archived two-period construction}
\label{supp:dynamic}

A one-horizon budget comparison omits one of the strongest reasons to expand: external industry can enlarge future resources. A minimal two-period construction makes this effect explicit while also allowing local investment to be productive. Let $B_1=10$, $z_1=B_1-p_1x$, and define
\begin{align}
 C_1(x)&=C_{01}+a_1\ln(1+z_1)+r_1x,\nonumber\\
 B_2(x)&=B_{20}+g_Lz_1+g_Xx,\label{eq:futurebudget}\\
 W(x)&=\ln[1+C_1(x)]+\beta V_2[B_2(x)],\quad 0\leq x\leq B_1/p_1.
 \label{eq:dynamic}
\end{align}
Here $B_1>0$, $B_{20},C_{01},C_{02}\geq0$, $a_1,a_2,p_1,p_2,r_1,r_2>0$, and $0\leq\beta\leq1$. The parameters $g_L,g_X\geq0$ convert current allocations to next-period resources. Both directions receive a future-production channel. All represented downstream benefits of first-period projects are captured through $B_2$; this is not a full persistent-capital or interstellar replication model.

The continuation value is the second-period pure-computation problem,
\begin{equation}
 V_2(B)=\max_{\substack{z_2,x_2\geq0\\z_2+p_2x_2\leq B}}\ln[1+C_{02}+a_2\ln(1+z_2)+r_2x_2].
 \label{eq:continuation}
\end{equation}
For $r_2>0$,
\begin{equation}
 z_2^*(B)=\min\{B,\max(a_2p_2/r_2-1,0)\},\qquad
 x_2^*(B)=\frac{B-z_2^*(B)}{p_2}.
 \label{eq:stage2}
\end{equation}
Its marginal value $V_2'(B)$ is the current best computational marginal return divided by $1+C_2^*(B)$: $a_2/(1+B)$ on the all-local branch, or $r_2/p_2$ on the mixed/external branch. The expressions agree at the switching budget. $V_2$ is increasing and concave.

The archived specialization of \cref{prop:dynamic} is
\begin{equation}
 \frac{r_1-p_1a_1/(1+B_1)}{1+C_{01}+a_1\ln(1+B_1)}
 +\beta V_2'(B_{20}+g_LB_1)(g_X-p_1g_L)\leq0.
 \label{eq:dynamiccriterion}
\end{equation}
The derivative proof in the main text applies because both current and
continuation objectives here are concave.


The first term is the original within-period comparison. The second values the \emph{difference} between external and local production of future resources. Counting $g_X$ while setting $g_L$ to zero without justification would bias the test against inward development; omitting $g_X$ would bias it in favour. This extension is a structural check, not a calibrated decision rule for actual civilizations.

For $p_1=p_2=r_1=r_2=1$, $a_1=a_2=100$, $C_{01}=C_{02}=B_{20}=1$, $g_L=0.2$, and $\beta=0.9$, present-only investment is all inward. The future-adjusted corner changes at $g_X=\DynamicThreshold$. At $g_X=0.2$ both periods remain all inward. At $g_X=2$, optimal first-period outward investment is $x=\DynamicOutward$, leaving $z_1=6.2441$ locally and producing $B_2=9.7607$. All second-period investment is still local. Thus the optimal path can include \emph{external investment undertaken to finance later inward activity}. This differs from claiming that external expansion is never useful. \Cref{fig:dynamic} gives the full 201-setting construction.

\begin{figure}[tbp]
\centering\includegraphics[width=.95\textwidth]{fig4.pdf}
\caption{Two-period resource feedback. Local investment generates future resources at $g_L=0.2$ per unit. Increasing the external yield $g_X$ eventually overturns the all-inward present-only optimum. Parameters are normalized; the transition is an exact derivative boundary checked against a separate three-variable numerical optimization.}
\label{fig:dynamic}
\end{figure}

The converse reversal is also possible. Keep
$B_1=10$, $p_1=p_2=r_1=r_2=1$ and $\beta=0.9$, but set
$a_1=5$, $a_2=100$, $C_{01}=C_{02}=1$, $B_{20}=0$,
$g_L=0.1$ and $g_X=0$. The static term in
\cref{eq:dynamiccriterion} is $\ReverseStaticTerm$, and the
present-only optimum is $x=\ReverseStaticX$. The continuation term is
$\ReverseFutureTerm$, making the total $\ReverseTotalTerm$.
Concavity therefore makes $x=\ReverseDynamicX$ the two-period optimum.
Local investment, rather than external investment, produces the more
valuable continuation budget in this example. Supplementary Material
section~S3 reports scalar and joint-optimization checks of both cases.

The two reversals are retained as finite-horizon examples. The infinite-horizon
benchmark in \cref{prop:infinite} now supplies a separate verification theorem,
with its own production and growth assumptions. The nonlinear examples here
are exact special cases of \cref{prop:dynamic}, not of the linear benchmark.

\section{Additional identities and new verification files}
\label{supp:newchecks}
The workload ratios used in the main text are, explicitly,
\begin{equation}
 \frac{D_k(d+1)}{D_k(d)}=b+\frac1{W_d}\longrightarrow b,
 \label{eq:recursion-cost}
\end{equation}
\begin{equation}
 F(C_{d+1})-F(C_d)=\ln\frac{1+c(bW_d+1)}{1+cW_d}\longrightarrow\ln b.
 \label{eq:recursion-value}
\end{equation}
Their proofs divide by $W_d\to\infty$. A nonconcave example retained from
v7.3 is
\begin{equation}
 U(x)=\ln\{1+2\ln(2-x)+0.1x\}+3x^2,\quad 0\leq x\leq1,
 \label{eq:nonconcave}
\end{equation}
with $U'(0)\simeq-0.3772$, $U(0)\simeq0.8697$, and
$U(1)\simeq3.0953$. It violates the concavity assumption, not the corner theorem.

Each new script in \texttt{code/validate\_v8\_*.py} is self-contained and
prints explicit pass and fail counts. The accompanying README maps files
to propositions and distinguishes assertions from independent randomized
cases; each script records its seed and tolerances. The new no-replication classification reads
\texttt{archive/v7\_3/data/allocation/design.npz} without changing it.
It solves observation on the restricted face, then evaluates the two
replication reduced derivatives; 64 of the 1,024 cases are additionally
compared with full constrained optimization. General restricted-face tests
use 320 separate randomized instances with optional backups. Thus the full
optimization check is sampled, whereas every frozen row receives the exact
face criterion.

The original archive scripts have been rerun in a working copy, not in the
immutable archive. Their outputs are included under
\texttt{results/archive\_rerun/}. No figures or archived classification
counts have been revised. Mechanical source and PDF checks are recorded in
\texttt{results/assembly\_audit.json} and
\texttt{results/final\_audit.json}; physical-parameter sources are documented
in the comments of \texttt{code/validate\_v8\_7\_physical.py}.
The finite-efficiency tests are in \texttt{code/validate\_v8\_8\_ceiling.py}
and \texttt{results/v8\_8\_ceiling.json}.
